\section{Introducción}
    Una prueba de hipótesis permite evaluar una afirmación sobre una población 
    a partir de una muestra mediante una hipótesis nula $H_0$ y una alternativa
    $H_1$, la decisión se basa en comparar el valor $p$ con un nivel de
    significancia $\alpha$ fijado previamente, sin embargo, un valor $p$ no 
    indica por sí solo la magnitud del efecto ni permite determinar si la 
    estructura de los datos es adecuada para la conclusión \cite{montgomery2020design},
    \cite{walpole2012probabilidad}, \cite{wasserstein2016asa}.

    \vspace{10pt}

    Esta limitación puede hacerse evidente cuando se combinan observaciones de
    diferentes subpoblaciones. El cuarteto de Anscombe muestra cómo conjuntos con
    estadísticas similares pueden presentar estructuras muy distintas al
    mostrarse de manera gráfica \cite{anscombe1973graphs}. De forma similar, la
    agregación de grupos con comportamientos diferentes puede ocultar o incluso
    modificar una asociación, como ocurre en la paradoja de Simpson
    \cite{kievit2013simpson}, \cite{simpson1951interpretation}, por ello, la
    visualización constituye una herramienta de diagnóstico y no solo un recurso
    para dar a conocer los resultados \cite{few2012show}, 
    \cite{castro2026visualizacion}.

    \vspace{10pt}

    Este informe reproduce este escenario mediante un conjunto simulado de 300
    clientes de los sectores Residencial, Comercial e Industrial, cuyas escalas
    de consumo son diferentes. Se aplican pruebas de hipótesis para comparar
    medias y analizar la relación entre temperatura y consumo, verificando
    previamente sus supuestos. La simulación permite además contrastar los
    resultados obtenidos con las relaciones definidas en los datos.

    \vspace{10pt}

    El análisis se desarrolla en Python con \textit{*SciPy*} y se replica en R,
    acompañando cada prueba con su respectiva visualización. El propósito es
    determinar cuándo un valor $p$ no significativo realmente indica ausencia de
    relación y cuándo una relación puede estar oculta por la estructura de los
    datos, para ello, se contrastan cinco hipótesis, se verifican los supuestos,
    se generan visualizaciones estáticas e interactivas y se interpretan los
    resultados desde una perspectiva estadística y práctica.